% ===================================================================
%  TABLEAU N° 10 — La Mer. — Le Port.
%  Feuillets 62 a 66 du fac-simile = folios 60 a 64.
%  Correspondance : folio imprime = numero de feuillet - 2.
%
%  COQUILLE DE NUMEROTATION (§7 de la consigne) — a partir du feuillet
%  63 (folio 61), le compositeur saute le numero 8 : l'imprime va de
%  « 7 --- » directement a « 9.--- », puis continue 10, 11, ... jusqu'a
%  « 26.--- » (25 alineas au total, comme dans l'ido, mais numerotes
%  1-7 puis 9-26 au lieu de 1-25). Le NUMERO IMPRIME EST CONSERVE TEL
%  QUEL dans le texte relevé (« 9.--- La langue de terre... », etc.).
%  En revanche, cette coquille decale d'un cran l'appariement avec
%  l'ido pour tout le reste du tableau : a partir de ce point, les
%  cles « %%K » portent donc le numero REEL d'alinea (celui qui
%  correspond au meme contenu que texto/io/19-tabelo-10.tex), et non
%  le numero fautif imprime — le bloc « 9.--- » porte la cle t10-08,
%  « 10.--- » porte t10-09, ... « 26.--- » porte t10-25. Verifie alinea
%  par alinea par comparaison de contenu avec l'ido (§6 : « une cle
%  fausse fait mentir la mise en regard »).
%
%  Le francais ne porte aucun intertitre interne (a la difference de
%  l'ido qui intercale « La mikra navi. », « La komercala portuo. »,
%  « La Arsenalo e la doki. », « La Piloto. », « La Remo-plezuro. »,
%  « La Klubeyo. », « Generala vidajo di la litori. ») : ecart REEL
%  entre les deux editions, non une lacune du releve.
%
%  Autres coquilles conservees telles quelles : « 7 --- » (sans point
%  apres le chiffre) ; « ballots 61) » (parenthese ouvrante absente,
%  meme defaut que « buvard 41) » au tableau 1) ; « passagers de
%  premiere classe. . » (point double, feuillet 64) ; le renvoi (82)
%  de « jeunes enseignes », a peine lisible sur le fac-similé, recoupe
%  par « sublietnanti (82) » dans texto/io/19-tabelo-10.tex.
%
%  ATTENTION — LE FAC-SIMILE FRANCAIS EST UNE PRISE DE VUE, contraste
%  inegal d'une page a l'autre : chaque page a ete verifiee a
%  l'agrandissement (2600 px) avant releve, et l'espace avant les
%  renvois a ete controlee ligne a ligne (elle varie meme au sein
%  d'une page : « croiseurs(6) » sans espace mais « transports (12) »
%  avec, par exemple).
% ===================================================================

% --- Feuillet 62 (folio 60) -----------------------------------------
\begin{VUpage}[62]{60}
%%K t10-apar-1 apar
\VUsaut{4.0mm}
\VUtitre{15.0pt}{60}{TABLEAU N\textsuperscript{o} 10}
\VUsaut{1.6mm}
\VUtitre{11.4pt}{0}{La Mer. --- Le Port.}
\VUsaut{3.0mm}

%%K t10-01-1 p
\VUblancAlinea
1. --- En été, tandis que les uns vont chercher dans\nl
la montagne le calme et la fraîcheur, les autres préfè\cc
rent se rendre au bord de la mer. C'est ce que nous\nl
avons fait l'année passée, car nous aimons beaucoup\nl
\VUgras{l'Océan} \textsuperscript{(1)}.

%%K t10-02-1 p
\VUblancAlinea
2. --- Pour ma part, je prends un vif plaisir à\nl
contempler cette immense nappe d'eau qui s'étend\nl
jusqu'à l'\VUgras{horizon} \textsuperscript{(2)}, et à voir partir pour une longue\nl
\VUgras{traversée} les énormes \VUgras{bateaux à vapeur} \textsuperscript{(3)} où\nl
s'embarquent les voyageurs qui vont en Amérique.

%%K t10-03-1 p
\VUblancAlinea
3. --- Aussi mon père, qui connaît mes goûts, choi\cc
sit-il toujours, pour passer les vacances, une petite\nl
plage bien tranquille, située près de l'un de nos\nl
grands \VUgras{ports de guerre.}

%%K t10-03-2 p
Pendant notre dernier séjour, l'\VUgras{escadre} était\nl
venue jeter l'ancre derrière l'énorme \VUgras{digue} \textsuperscript{(4)} qui\nl
ferme et protège le \VUgras{port militaire} \textsuperscript{(5)}, et j'ai pu\nl
admirer à loisir les différents \VUgras{navires de guerre.}\nl
depuis le petit \VUgras{sous-marin} \textsuperscript{(10)} qui plonge et disparaît\nl
sous les eaux, jusqu'à l'énorme \VUgras{cuirassé} \textsuperscript{(9)} qui,\nl
jusqu'à présent, ne redoutait guère que les attaques\nl
du \VUgras{torpilleur} \textsuperscript{(8)}.

%%K t10-04-1 p
\VUblancAlinea
4. --- Celui-ci savait déjà fort habilement découvrir\nl
le point faible de l'épaisse cuirasse métallique pour y\nl
planter la dangereuse \VUgras{torpille} dont l'explosion cause\nl
la perte du bâtiment, mais il était cependant moins\nl
redoutable que le sous-marin, car les contre-torpilleurs\nl
lui donnaient facilement la chasse.

%%K t10-05-1 p
\VUblancAlinea
5. --- J'ai pu voir aussi les rapides \VUgras{croiseurs}\textsuperscript{(6)}\nl
cuirassés, presque aussi invulnérables que le véritable\nl
cuirassé, plusieurs \VUgras{transports} \textsuperscript{(12)}, trois ou quatre\nl
\VUgras{canonnières}\textsuperscript{(7)} et enfin une \VUgras{frégate}\textsuperscript{(11)}. Le spectacle
\parplein
\end{VUpage}

% --- Feuillet 63 (folio 61) -----------------------------------------
\begin{VUpage}[63]{61}
%%K t10-05-1 p suite
\VUcontinue
de cette flotte, lorsqu'elle appareille pour lever\nl
l'ancre, est des plus intéressants.

%%K t10-06-1 p
\VUblancAlinea
6. --- Quelle différence de construction entre ces\nl
énormes masses d'acier et les élégants \VUgras{trois-mâts}\nl
ou les fins \VUgras{voiliers} \textsuperscript{(13)}, employés encore dans la\nl
marine marchande, que je voyais entrer dans le port,\nl
toutes voiles dehors, lorsque j'allais me promener sur\nl
la \VUgras{jetée} \textsuperscript{(14)}. Ceux-ci, il est vrai, perdaient beaucoup\nl
de leurs avantages lorsque le vent était contraire,\nl
qu'il fallait plier toutes leurs voiles pour entrer dans\nl
le bassin et qu'un \VUgras{remorqueur} \textsuperscript{(16)} était obligé d'aller\nl
les prendre pour les amener à quai comme de simples\nl
\VUgras{chalands} \textsuperscript{(17)}.

%%K t10-07-1 p
\VUblancAlinea
7 --- Ce qui fait l'intérêt du port dont je parle, c'est\nl
qu'il ne comprend pas seulement un port de guerre,\nl
\VUgras{créé} artificiellement à l'aide de travaux gigantesques,\nl
mais encore un merveilleux \VUgras{port de commerce,}\nl
situé à l'\VUgras{embouchure} d'un grand fleuve.

%%K t10-08-1 p
\VUblancAlinea
9.--- La langue de terre qui sépare les deux bassins\nl
est fortifiée et défendue par une série de \VUgras{batteries}\nl
et de \VUgras{bastions} qui s'avancent dans la mer. Il faut\nl
une autorisation spéciale pour se promener sur les\nl
\VUgras{remparts} \textsuperscript{(15)}. Je l'avais obtenue sans peine, par\nl
mon oncle, qui est ingénieur des \VUgras{chantiers de}\nl
\VUgras{construction} \textsuperscript{(18)}, et je suis allé souvent visiter les\nl
bateaux que l'on construit sous de vastes hangars.

%%K t10-09-1 p
\VUblancAlinea
10.--- J'ai même assisté au lancement d'un cuirassé\nl
et je me souviens du moment d'émotion que toute la\nl
foule éprouva lorsque l'énorme masse, abandonnée à\nl
elle-même, se mit à glisser sur son ber et vint flotter\nl
sur l'eau du fleuve.

%%K t10-10-1 p
\VUblancAlinea
11. --- Pour aller aux chantiers, on traverse le ter\cc
rain de \VUgras{manoeuvres} \textsuperscript{(19)}, où les troupes de la gar\cc
nison viennent chaque jour faire l'exercice, et on\nl
passe sur une \VUgras{passerelle} \textsuperscript{(20)} en fer, qui relie le\nl
Champ-de-Mars à l'\VUgras{arsenal} \textsuperscript{(21)}.
\end{VUpage}

% --- Feuillet 64 (folio 62) -----------------------------------------
\begin{VUpage}[64]{62}
%%K t10-11-1 p
\VUblancAlinea
12. --- J'ai visité, avec mon oncle, ce grand établis\cc
sement plein de \VUgras{canons} \textsuperscript{(22)}, de \VUgras{boulets} \textsuperscript{(23)} et d'\VUgras{obus.}\nl
J'ai vu aussi l'\VUgras{usine métallurgique} \textsuperscript{(24)} où sont\nl
fabriquées les \VUgras{chaudières} \textsuperscript{(25)} des bateaux à vapeur.

%%K t10-12-1 p
\VUblancAlinea
13. --- Tout près de là se trouvent les \VUgras{bassins} \textsuperscript{(26)}\nl
(formes de radoub) et les très curieux \VUgras{docks flot}\cc
\VUgras{tants} \textsuperscript{(27)} dans lesquels j'ai pu voir de près, sur un\nl
navire en réparation, l'\VUgras{hélice} \textsuperscript{(28)}, la \VUgras{quille} \textsuperscript{(29)} et le\nl
\VUgras{gouvernail} \textsuperscript{(30)} d'un bateau.

%%K t10-13-1 p
\VUblancAlinea
14. --- Le port de commerce est tout aussi intéressant\nl
à étudier que le port de guerre, et j'ai passé bien des\nl
heures à flâner sur les quais où sont amarrés les\nl
\VUgras{bateaux à aubes} \textsuperscript{(31)} qui remontent le fleuve, et à\nl
regarder fonctionner la \VUgras{grue à mâter} \textsuperscript{(32)}, qui sou\cc
lève comme des plumes des fardeaux énormes.

%%K t10-14-1 p
\VUblancAlinea
15. --- J'admirais les \VUgras{transatlantiques} \textsuperscript{(34)} sortant\nl
de la rade, \VUgras{pavillon} \textsuperscript{(35)} au vent, conduits par un\nl
habile \VUgras{pilote} \textsuperscript{(36)} qui savait merveilleusement suivre\nl
le \VUgras{chenal} marqué par les \VUgras{bouées} \textsuperscript{(40)} et les \VUgras{balises,}\nl
en évitant les \VUgras{gabares} \textsuperscript{(41)} qui se dirigent pénible\cc
ment avec leur unique voile carrée. Je distinguais\nl
à l'\VUgras{avant} \textsuperscript{(37)} (proue) les \VUgras{cabines} \textsuperscript{(39)} de deuxième\nl
classe, et, à l'\VUgras{arrière} \textsuperscript{(38)} (poupe), les salons pour les\nl
passagers de première classe. .

%%K t10-15-1 p
\VUblancAlinea
16. --- Parfois aussi j'assistais au chargement d'un\nl
\VUgras{cargo-boat} \textsuperscript{(42)} où venaient s'entasser les marchan\cc
dises de l'\VUgras{entrepôt} \textsuperscript{(43)} et de la \VUgras{gare maritime} \textsuperscript{(44)},\nl
amenées par les nombreux trains de la \VUgras{ligne des}\nl
\VUgras{quais} \textsuperscript{(45)}.

%%K t10-16-1 p
\VUblancAlinea
17. --- J'ai souvent canoté dans le \VUgras{bassin à flot} \textsuperscript{(53)}\nl
où viennent se réfugier les petites \VUgras{embarcations} \textsuperscript{(46)},\nl
et c'était une joie pour moi de manier l'\VUgras{aviron} \textsuperscript{(47)}.\nl
Je montais sur le \VUgras{pont} \textsuperscript{(48)} d'une barque à \VUgras{voile} \textsuperscript{(49)},\nl
je grimpais au \VUgras{mât} \textsuperscript{(51)} et dans les \VUgras{vergues} \textsuperscript{(50)} à\nl
l'aide des \VUgras{cordages} \textsuperscript{(52)}, puis je reprenais ma prome\cc
nade en \VUgras{yole} \textsuperscript{(54)} au milieu des lourdes \VUgras{barques de}\nl
\VUgras{pêcheurs} \textsuperscript{(55)}, où sèchent les \VUgras{filets} \textsuperscript{(56)}.
\end{VUpage}

% --- Feuillet 65 (folio 63) -----------------------------------------
\begin{VUpage}[65]{63}
%%K t10-17-1 p
\VUblancAlinea
18. --- Sur le \VUgras{quai} \textsuperscript{(58)} défilaient sous mes yeux les\nl
\VUgras{marchandises} \textsuperscript{(59)} les plus diverses, contenues dans\nl
des \VUgras{caisses} \textsuperscript{(60)} ou dans des \VUgras{ballots} \textsuperscript{61)}. Elles for\cc
maient la \VUgras{cargaison} \textsuperscript{(63)} des péniches et des \VUgras{grues}\textsuperscript{(62)}\nl
puissantes les enlevaient et les déposaient sur des\nl
\VUgras{camions} \textsuperscript{(64)} pendant que les \VUgras{camionneurs} faisaient\nl
leur déclaration et payaient les droits au \VUgras{bureau de}\nl
\VUgras{la douane} \textsuperscript{(65)}.

%%K t10-18-1 p
\VUblancAlinea
19. --- Enfin, lorsque j'arrivais au \VUgras{pont-tour}\cc
\VUgras{nant} \textsuperscript{(66)} ou à l'\VUgras{écluse} \textsuperscript{(69)}, en passant devant la\nl
\VUgras{drague} \textsuperscript{(68)} qui nettoie le \VUgras{canal} \textsuperscript{(67)}, je débarquais sur\nl
la jetée à côté du \VUgras{sémaphore} \textsuperscript{(70)}.

%%K t10-19-1 p
\VUblancAlinea
20. --- C'est de l'autre côté de cette jetée que vien\cc
nent accoster les \VUgras{chaloupes} \textsuperscript{(71)} qui amènent à terre\nl
les officiers de l'escadre. C'est un spectacle curieux de\nl
voir les matelots lever tous ensemble leurs avirons,\nl
tandis que l'embarcation, portée par la \VUgras{vague} \textsuperscript{(72)},\nl
vient aborder sans peine à l'échelle du débarcadère.\nl
C'est à cet endroit que l'on peut le mieux observer la\nl
\VUgras{marée} et le mouvement du \VUgras{flux} et du \VUgras{reflux} sur la\nl
\VUgras{plage} \textsuperscript{(73)}.

%%K t10-20-1 p
\VUblancAlinea
21. --- Pour rentrer chez mes parents, je prenais le\nl
\VUgras{tramway électrique} \textsuperscript{(74)}, dont la \VUgras{halte} \textsuperscript{(75)} est située\nl
au \VUgras{poteau} placé devant le cercle des officiers.

%%K t10-20-2 p
Sur le \VUgras{balcon} \textsuperscript{(76)} du cercle, décoré d'\VUgras{ancres}\textsuperscript{(77)},\nl
de canons et de boulets, j'apercevais souvent une\nl
brillante réunion d'officiers de marine : \VUgras{lieutenants}\nl
\VUgras{de vaisseau} \textsuperscript{(78)} portant l'épée, \VUgras{capitaines d'artil}\cc
\VUgras{lerie} \textsuperscript{(79)} et d'\VUgras{infanterie} \textsuperscript{(80)} de marine avec leur\nl
\VUgras{sabre.} Derrière eux se tenait un \VUgras{matelot} \textsuperscript{(81)} qui\nl
leur apportait une \VUgras{lunette d'approche,} lorsqu'ils\nl
voulaient distinguer ce qui se passait au loin.

%%K t10-21-1 p
\VUblancAlinea
22. --- Je voyais aussi de jeunes \VUgras{enseignes} \textsuperscript{(82)} pre\cc
nant les ordres de leur \VUgras{capitaine de vaisseau} \textsuperscript{(83)}\nl
ou de l'\VUgras{amiral} \textsuperscript{(84)}, tandis qu'un \VUgras{quartier-maître}\textsuperscript{(85)}\nl
attendait pour les porter à bord.
\end{VUpage}

% --- Feuillet 66 (folio 64) -----------------------------------------
\begin{VUpage}[66]{64}
%%K t10-22-1 p
\VUblancAlinea
23. --- De ce balcon, la vue est magnifique : on\nl
domine toute la plage avec ses \VUgras{tentes} \textsuperscript{(86)} et ses\nl
\VUgras{cabines} \textsuperscript{(87)}, et l'on voit, à l'heure du bain, les bai\cc
\VUgras{gneurs} \textsuperscript{(88)}, qui s'ébattent joyeusement en face du\nl
\VUgras{casino} \textsuperscript{(89)}.

%%K t10-23-1 p
\VUblancAlinea
24. --- A l'extrémité de la plage, la côte forme une\nl
petite \VUgras{presqu'île}\textsuperscript{(91)} de \VUgras{rochers}, où la \VUgras{lame}\textsuperscript{(92)} vient\nl
se briser, et sur laquelle est bâti un \VUgras{phare} \textsuperscript{(93)}, qui\nl
indique, la nuit, la route aux navigateurs. Cette\nl
presqu'île n'est reliée à la terre que par un \VUgras{isthme}\textsuperscript{(90)}\nl
très étroit.

%%K t10-24-1 p
\VUblancAlinea
25. --- La \VUgras{falaise} \textsuperscript{(94)} se prolonge, déchiquetée par\nl
la mer, qui y dessine capricieusement une série de\nl
\VUgras{baies} \textsuperscript{(95)} et de \VUgras{promontoires} (caps), qui donnent\nl
à la côte un aspect des plus pittoresques.

%%K t10-24-2 p
Sur le \VUgras{plateau} \textsuperscript{(96)}, qui domine le \VUgras{détroit} \textsuperscript{(98)}, il y\nl
a un \VUgras{fort} \textsuperscript{(97)} très important et quelques villas.

%%K t10-25-1 p
\VUblancAlinea
26.--- Ce détroit sépare du continent une \VUgras{île} \textsuperscript{(99)} très\nl
dangereuse, bordée d'\VUgras{écueils} \textsuperscript{(100)} et de \VUgras{récifs}, où des\nl
barques et même de grands navires viennent parfois\nl
s'échouer. Malgré la promptitude des secours et\nl
l'arrivée rapide des \VUgras{canots de sauvetage}, ces\nl
\VUgras{naufrages} sont terribles, et, par les fortes \VUgras{tempê}\cc
\VUgras{tes} d'équinoxe, plusieurs navires ont été perdus corps\nl
et biens.

%%K t10-25-2 p
En dépit de ce voisinage, le port est très fréquenté\nl
par les marins.

\VUsaut{6.0mm}
\VUfilet{22mm}
\end{VUpage}
